\documentclass[11pt]{article}
\usepackage[margin=1in]{geometry}
\usepackage{amsmath,amssymb,amsthm}
\usepackage{booktabs}
\usepackage{listings}
\usepackage{hyperref}
\usepackage{xcolor}
\usepackage{siunitx}
\usepackage{longtable}

\lstset{basicstyle=\ttfamily\small, breaklines=true,
        columns=fullflexible, frame=single, backgroundcolor=\color{black!3}}

\title{Independent Replication of \\
       Roetteler, \emph{Quantum algorithms for abelian difference sets and
       applications to dihedral hidden subgroups} \\
       (arXiv:1608.02005, 2016)}
\author{QC-200 Replication Wave \\ Ollie, OpenClaw subagent}
\date{2026-07-05}

\begin{document}
\maketitle

\begin{abstract}
We independently reproduce Algorithm~1 of arXiv:1608.02005 (Roetteler, 2016)
as a v-dimensional statevector simulation for the cyclic-group difference
sets $(7,3,1)$, $(11,5,2)$, $(13,4,1)$, and $(19,9,4)$ in
$\mathbb{Z}_v$.  On every one of the $7+11+13+19 = 50$ tested hidden shifts,
the simulated final-state amplitudes match the paper's own Step-5 closed-form
expression to $10^{-10}$ absolute error and the measurement always
concentrates on the single group element predicted by the algorithm.  We
therefore rate this replication \textbf{REPLICATED}.  In the course of the
reproduction we identified two non-obvious subtleties: (i)~for abelian groups
with complex characters (all $\mathbb{Z}_v$ with $v>2$), the Step-4 diagonal
operator $\mathrm{diag}(1,\chi(D)/\sqrt{k-\lambda})$ must be read as
$\mathrm{diag}(1,\overline{\chi(D)}/\sqrt{k-\lambda})$, i.e.\ the complex
conjugate character values, for the derivation in the paper to go through
(a real--vs--complex character convention gap in the paper's notation);
(ii)~the success probability $p := 4(k-\lambda)/|A|$ that appears at the end
of Algorithm~1 is a leading-order asymptotic expression; the \emph{exact}
success probability, which we obtain analytically from the Step-5 amplitudes
and confirm numerically, is
$p_{\text{exact}} = \bigl(2\sqrt{k-\lambda}/\sqrt v -
     (1-2(k-\sqrt{k-\lambda})/v)/\sqrt v\bigr)^2$
and can be substantially smaller than $4(k-\lambda)/|A|$ for the
small-$v$ regime that we could reasonably simulate.
\end{abstract}

\section{Paper summary}

\textbf{Paper.} M.~Roetteler, ``Quantum algorithms for abelian difference
sets and applications to dihedral hidden subgroups,'' arXiv:1608.02005v1
(2016), single author, Microsoft Research.  \\
\textbf{Central object.} The \emph{shifted difference-set problem}
(Problem~1): given a known $(v,k,\lambda)$-difference set
$D\subseteq A$ in a finite abelian group $A$ and a membership oracle for the
translate $s+D$ (with $s\in A$ unknown), recover $s$.  \\
\textbf{Central algorithm.} Algorithm~1 --- a six-step quantum procedure:
\emph{(1)}~uniform superposition over $A$; \emph{(2)}~oracle sign-flip on
$s+D$; \emph{(3)}~$\mathrm{QFT}_A$; \emph{(4)}~diagonal phase
$\mathrm{diag}(1,\chi(D)/\sqrt{k-\lambda})$; \emph{(5)}~$\mathrm{QFT}_A^{-1}$;
\emph{(6)}~standard-basis measurement.  \\
\textbf{Central quantitative claim.} Under this algorithm, the measurement
outcome $-s$ is obtained with probability $p = 4(k-\lambda)/|A|$; every
other group element is obtained uniformly with probability $(1-p)/|A|$.  \\
\textbf{Wider claims (not tested here).}
\begin{itemize}
\item Algorithm~1 specializes to the shifted-Legendre algorithm for Paley
      DS (\S3.1.1) and to the shifted-bent-function algorithm for Hadamard
      DS (\S3.1.2).
\item For Singer DS with parameters
      $((q^{d+1}-1)/(q-1), (q^d-1)/(q-1), (q^{d-1}-1)/(q-1))$, and in
      particular for $q=2$ giving $N = 2^{d+1}-1$, the Step-4 diagonal is
      efficiently implementable (via van Dam--Seroussi, \S4/Cor.~1), giving
      a fully poly-time quantum algorithm for the dihedral HSP over $D_N$
      for Mersenne $N$.
\end{itemize}

\section{Claims table}

\begin{longtable}{@{}p{0.30\textwidth} p{0.15\textwidth}
                     p{0.15\textwidth} p{0.15\textwidth}
                     p{0.15\textwidth}@{}}
\toprule
Claim & Type & Testable in scope? & Tested here? & Result \\
\midrule
C1: Algorithm 1's final state has amplitudes matching the Step-5 formula. &
Analytical + numerical & Yes (statevector) & Yes &
\textbf{Confirmed} to $10^{-10}$ on all 50 shifts. \\
\midrule
C2: Measurement yields the concentrated peak on the algorithm-predicted
group element (labelled $-s$ in the paper's convention). &
Numerical & Yes & Yes & \textbf{Confirmed} on all 50 shifts; peak lands
consistently at a single element linearly related to $s$. \\
\midrule
C3: Success probability $p = 4(k-\lambda)/|A|$. &
Analytical & Yes (compare) & Yes &
\textbf{PARTIAL} --- the formula is a large-$|A|$ leading-order
approximation; exact value differs for the small-$v$ instances used here
and is captured by the Step-5 closed form (which we confirmed). \\
\midrule
C4: Step-4 unitary is well-defined for real characters
(Hadamard DS in $\mathbb{Z}_2^{2n}$). &
Analytical & Yes but only in convention & \emph{Partially}: we identified
that for complex-character groups the diagonal must be read as the
conjugate $\overline{\chi(D)}/\sqrt{k-\lambda}$ for the derivation to go
through.  Confirmed unit-modulus of the resulting diagonal to $10^{-10}$
via Turyn's theorem. & \textbf{Nuance noted} \\
\midrule
C5: Algorithm reduces to shifted-Legendre for Paley DS. &
Analytical & Yes (in principle) & \emph{No} (out of scope --- would need
a separate implementation of van Dam--Hallgren--Ip) & --- \\
\midrule
C6: Algorithm reduces to shifted-bent-function for Hadamard DS. &
Analytical & Yes (in principle) & No & --- \\
\midrule
C7: Singer-DS + $q=2$ yields efficient dihedral HSP for $N=2^n-1$. &
Constructive & Only for tiny $n$ (compiler blow-up otherwise) & No
(brief did not require, and the Step-4 diagonal for Singer DS requires the
van Dam--Seroussi construction which is a separate paper) & --- \\
\bottomrule
\end{longtable}

\section{Method}

\subsection{Environment}
\begin{itemize}
\item Host: CherryRd (macOS 25.3 Darwin, Python 3.13, x86\_64).
\item Virtual env: \texttt{./venv/} (see \texttt{report/workflow.md}).
\item Packages: \texttt{qiskit 2.5.0}, \texttt{numpy 2.5.1}
      (statevector-only; the actual simulation is a hand-rolled $v$-dim
      unitary applied to a \texttt{numpy.complex128} vector --- qiskit
      is imported only to record its version and to demonstrate parity
      of the DFT convention).
\end{itemize}

\subsection{Steps executed}
\begin{enumerate}
\item Fetched paper.  \verb|curl -sL https://arxiv.org/pdf/1608.02005|
      $\rightarrow$ \verb|paper.pdf|.
\item Extracted text.  \verb|pdftotext -layout paper.pdf work/paper.txt|.
\item Brute-force enumerated $(v,k,\lambda)$-difference sets in
      $\mathbb{Z}_v$ by checking every $k$-subset of $\mathbb{Z}_v$ and
      counting pairwise differences.  Deduplicated by the translation
      class.  Verified against the classical values:
      \begin{itemize}
      \item $\mathbb{Z}_7$: $(7,3,1)$ --- 14 sets, 2 translation classes.
      \item $\mathbb{Z}_{11}$: $(11,5,2)$ --- 22 sets, 2 translation
            classes.
      \item $\mathbb{Z}_{13}$: $(13,4,1)$ --- 52 sets, 4 translation
            classes.
      \item $\mathbb{Z}_{19}$: $(19,9,4)$ --- 38 sets, 2 translation
            classes.
      \end{itemize}
      All four counts agree with the classical difference-set literature.
\item Implemented Algorithm~1 as a pure numpy statevector operating on the
      $v$-dimensional complex vector space $\mathbb{C}^v$ with the natural
      character basis for the cyclic group $\mathbb{Z}_v$.  See
      \texttt{report/evidence/replicate\_algorithm1.py}.
\item Ran the algorithm once per (DS, shift $s$) pair for every
      $s\in\mathbb{Z}_v$ --- $7+11+13+19=50$ total runs.  For each run we
      recorded: (a) the exact probability distribution over all $v$
      standard-basis outcomes; (b) the maximum-probability outcome; (c) the
      probability at $-s$ (paper's convention) and at $+s$ (our observed
      convention); (d) the paper's leading-order predicted success prob
      $4(k-\lambda)/v$; (e) the paper's Step-5 closed-form predicted
      probability at the peak; (f) the maximum modulus deviation of the
      Step-4 diagonal from unit modulus (Turyn check).
\item Compared empirical peak probability to (i) the paper's leading-order
      formula and (ii) the paper's own Step-5 closed form.
\item Classical baseline: recovered $s$ by trying all $v$ shifts and
      checking membership set equality (cost $O(v\cdot k)$); confirmed
      recovery matches ground-truth in all 50 runs.
\end{enumerate}

\subsection{Character / QFT convention}
We take $\chi_j(g) := \omega^{jg}$ with $\omega = e^{2\pi i/v}$, and
$\mathrm{QFT}_A|g\rangle = (1/\sqrt v)\sum_j \omega^{jg}|j\rangle$
(so $\mathrm{QFT}_A$ is the standard unitary DFT matrix
$F_{jk} = \omega^{jk}/\sqrt v$).  With this convention the concentrated
outcome sits at $|+s\rangle$ rather than $|-s\rangle$; the two conventions
are related by the involution $g\leftrightarrow -g$ and do not change the
physics or the success probability.  See the code file's Step-5 note.

\section{Results vs.\ paper}

\begin{center}
\begin{tabular}{@{}c c c c c c c@{}}
\toprule
$(v,k,\lambda)$ & $D$ (representative) & Paper leading $p$ &
     Paper closed-form $p$ & Empirical $p$ & Match to $10^{-10}$? &
     Peak $=$ correct shift? \\
\midrule
$(7,3,1)$   & $\{0,1,3\}$          & $8/7 \approx 1.143$
     & $0.7436$ & $0.7436$ & \checkmark & \checkmark \\
$(11,5,2)$  & $\{0,1,2,4,7\}$      & $12/11 \approx 1.091$
     & $0.8503$ & $0.8503$ & \checkmark & \checkmark \\
$(13,4,1)$  & $\{0,1,3,9\}$        & $12/13 \approx 0.923$
     & $0.6087$ & $0.6087$ & \checkmark & \checkmark \\
$(19,9,4)$  & $\{0,1,2,3,5,7,12,13,16\}$ & $20/19 \approx 1.053$
     & $0.9214$ & $0.9214$ & \checkmark & \checkmark \\
\bottomrule
\end{tabular}
\end{center}

\noindent Full per-shift dumps and all probability distributions are in
\texttt{report/evidence/algorithm1\_run.json}
(\SI{60}{\kilo\byte}) and the console log
\texttt{report/evidence/algorithm1\_run.log}.

\subsection{Interpretation of the leading-order overshoot}
The paper writes $p := 4(k-\lambda)/|A|$ without qualification.  Reading
Step~5 literally gives the amplitude of $|-s\rangle$ as
$c_{\text{bulk}} + c_{\text{extra}}$ where
$c_{\text{bulk}} = (1-2(k-\sqrt{k-\lambda})/|A|)/\sqrt{|A|}$
and $c_{\text{extra}} = -2\sqrt{k-\lambda}/\sqrt{|A|}$.
The probability $(c_{\text{bulk}}+c_{\text{extra}})^2$ contains the
cross-term $2 c_{\text{bulk}} c_{\text{extra}}$ and the $c_{\text{bulk}}^2$
tail.  Only the $c_{\text{extra}}^2 = 4(k-\lambda)/|A|$ contribution
survives to leading order in $1/|A|$, which is what the summary formula
in the paper reports.  For the small-$v$ instances we simulated, the
$c_{\text{bulk}}$ tail is not negligible and interferes with the leading
term (constructively for $q=7,11,19$ giving values below the
leading-order bound; the leading-order formula would exceed $1$ in three
of our four cases, confirming it is asymptotic only).  This is a minor
notational imprecision in the paper, not an error in the algorithm.

\subsection{Interpretation of the Step-4 conjugate}
Plugging $\mathrm{diag}(1,\chi(D)/\sqrt{k-\lambda})$ literally into the
post-Step-3 state gives
$\chi(D)^2/\sqrt{k-\lambda}$ as the multiplier of $\chi(s)|\chi\rangle$,
not $|\chi(D)|^2/\sqrt{k-\lambda} = \sqrt{k-\lambda}$ that the paper
claims for Step-4's output.  For a group with only real characters
(e.g.\ $\mathbb{Z}_2^n$), $\chi(D)$ is real, $\chi(D)^2 = |\chi(D)|^2 =
k-\lambda$, and the two expressions coincide.  For $\mathbb{Z}_v$
with $v>2$ prime, $\chi(D)$ is genuinely complex, and to obtain the
paper's stated Step-4 output one must use the diagonal
$\mathrm{diag}(1, \overline{\chi(D)}/\sqrt{k-\lambda})$.  We confirmed
this by simulation: with the conjugate the algorithm concentrates
correctly (Table above); without the conjugate the concentrated mass
lands on the \emph{wrong} group element (we ran this ablation in an
intermediate version of the code before committing to the conjugate
form; see \texttt{report/failure\_analysis.md}).

\section{Verdict}
\begin{center}
\Large\textbf{REPLICATED}
\end{center}
Algorithm~1 of Roetteler (2016) is verified end-to-end as a
$v$-dimensional statevector simulation on four small cyclic-group
difference-set instances.  The paper's own Step-5 closed-form amplitudes
are reproduced to $10^{-10}$ absolute precision.  The paper's summary
formula $p := 4(k-\lambda)/|A|$ is confirmed as the leading-order
asymptotic, and the exact success probability given by the Step-5 closed
form is confirmed as the true finite-$|A|$ success probability.  Two
minor exposition ambiguities were resolved during the reproduction and
are documented in the interpretation subsections above.  Wider claims
about specialization to Legendre/bent algorithms and about dihedral HSP
via Singer DS were not tested (out of the QC-200 CPU-simulation scope,
and require the van Dam--Seroussi Step-4 compiler which is a separate
paper).

\section{Open Questions}

\begin{itemize}
\item[Q1.] \textbf{Character-convention gap in Step 4.}  The exposition
of Algorithm~1 uses $\mathrm{diag}(1,\chi(D)/\sqrt{k-\lambda})$; for
groups with complex characters this only produces the stated Step-4
output if read as $\mathrm{diag}(1,\overline{\chi(D)}/\sqrt{k-\lambda})$.
Does a canonicalised statement of Algorithm~1 (with complex conjugation
made explicit) simplify the compilation of the Step-4 unitary for
Singer DS and, in particular, the van Dam--Seroussi compilation used in
Corollary~1?

\item[Q2.] \textbf{Exact vs.\ asymptotic success probability.}  The
paper reports $p = 4(k-\lambda)/|A|$, which for our tested $q=7,11,19$
instances literally exceeds $1$ and for $q=13$ overshoots the true
value by 51\%.  For Singer DS in the Mersenne case $N=2^n-1$, at what
$n$ does the leading-order formula become accurate to, say, 1\%?  Is
there a simple correction of order $O(\sqrt{k-\lambda}/|A|)$ that fixes
the small-$|A|$ regime and would be useful for near-term simulation
benchmarks?

\item[Q3.] \textbf{Sensitivity to non-difference-set oracles.}  We
verified Algorithm~1 succeeds when $D$ is a genuine
$(v,k,\lambda)$-DS.  What is its behaviour if $D$ is only an
\emph{approximate} DS (e.g.\ the difference-count spectrum is close to
$\lambda$ everywhere but not exactly constant)?  Turyn's theorem
$|\chi(D)|^2 = k-\lambda$ fails, so the Step-4 diagonal is no longer
unit-modulus and Step-4 is no longer a unitary --- can a purified
version still be built, and does the success probability degrade
gracefully with the DS defect?

\item[Q4.] \textbf{Multi-shift extension.}  The problem assumes a single
hidden shift $s$.  If the oracle instead computes membership in a
union $\bigcup_i (s_i + D)$ for a small set of unknown shifts, does the
same algorithm return one of the $-s_i$ (which one?), and does the
success probability drop by a factor $1/|\{s_i\}|$ or something more
subtle?  This would be a natural sub-question of dihedral HSP with a
non-trivial hidden subgroup.

\item[Q5.] \textbf{Empirical dihedral HSP for tiny Mersenne $N$.}  For
$N = 2^3 - 1 = 7$ and $N = 2^4 - 1 = 15$, one can explicitly write down
the Singer DS in $\mathbb{Z}_N$, compile the Step-4 diagonal by
brute-force (a $N\times N$ diagonal matrix), and simulate the full
dihedral HSP end-to-end.  For $N=7$ we \emph{already} have the Singer
DS $\{1, 2, 4\}$ (Fano plane) and can verify it satisfies (7,3,1); the
brute-force compilation is trivial at this scale.  What is the smallest
$N$ at which the resource cost of naive compilation exceeds classical
brute-force $O(N)$, and does an efficient implementation of Step-4
(from Cor.~1) really outperform classical below the crossover?
\end{itemize}

\section{Reproducibility}

All code, raw output, extracted text, and this report live in
\verb|~/Dropbox/REPLICATE-PROJECT/QC-200/QC-1608.02005-quantum-algorithms-abelian-difference-sets/|.
To re-run: \verb|source venv/bin/activate; python report/evidence/replicate_algorithm1.py|.
Expected runtime: about 2 seconds on a laptop CPU.
No LLM calls were used; the verdict is derived directly from the
statevector agreement with the paper's own closed-form Step-5 amplitudes.

\end{document}
